\documentclass[11pt]{article}

\usepackage[margin=1in]{geometry}
\usepackage{amsmath,amssymb,amsthm,mathtools}
\usepackage{graphicx}
\usepackage{booktabs}
\usepackage{caption}
\usepackage{subcaption}
\usepackage{hyperref}
\usepackage{float}
\usepackage{enumitem}
\usepackage{microtype}

\hypersetup{colorlinks=true,linkcolor=blue,citecolor=blue,urlcolor=blue}

% Notation
\newcommand{\R}{\mathbb{R}}
\newcommand{\E}{\mathbb{E}}
\newcommand{\1}{\mathbf{1}}
\newcommand{\DeltaM}{\Delta_m}
\newcommand{\Deltasparse}{\Delta_{m,s}}
\newcommand{\supp}{\operatorname{supp}}
\newcommand{\argmin}{\operatorname*{arg\,min}}
\newcommand{\argmax}{\operatorname*{arg\,max}}
\newcommand{\eps}{\varepsilon}
\newcommand{\Sstar}{S^\star}
\newcommand{\alphastar}{\alpha^\star_s}
\newcommand{\norm}[1]{\left\lVert #1\right\rVert}
\newcommand{\DeltaK}{\Delta_\kappa}
\newcommand{\DeltaF}{\Delta_f}
\newcommand{\DeltaLoss}{\Delta_L}

\title{Experimental Evaluation of Sparse Mixture-UCB Algorithms\\
on 2D Gaussian Generators}
\author{}
\date{\today}

\begin{document}
\maketitle

\begin{abstract}
We compare three Sparse Mixture-UCB algorithms---Fully-Corrective Sparse (Algorithm~2), Coherence-Structured Forced-Exploration (Algorithm~3), and the non-sparse Mixture-UCB baseline---on a toy problem with $m=8$ fixed 2D isotropic Gaussian generators. Two objectives are evaluated: the target-free RKE diversity objective and a squared-MMD objective with a known 3-component Gaussian-mixture target. The coherence-structured algorithm leverages the factorized least-squares structure of MMD to recover the oracle support via NNOMP-style residual-correlation scoring, then concentrates samples on that support. It achieves a 7--14$\times$ regret reduction over the next-best sparse method on MMD, while matching fully-corrective performance on RKE.
\end{abstract}

\section{Experimental Setup}

\subsection{Generators}
The experiment uses a fixed set of $m=8$ bivariate Gaussian generators,
\[
P_i = \mathcal{N}(\mu_i, \sigma_g^2 I_2), \qquad \sigma_g = 0.40,
\]
with means given in Table~\ref{tab:means}. All generators share the same isotropic covariance. Samples are drawn independently at each round.

\begin{table}[H]
\centering
\begin{tabular}{c|c}
\toprule
$i$ & $\mu_i$ \\
\midrule
1 & $(-3.0, -2.0)$ \\
2 & $(-2.0, \phantom{+}1.5)$ \\
3 & $(-0.8, -1.2)$ \\
4 & $(\phantom{+}0.0, \phantom{+}2.4)$ \\
5 & $(\phantom{+}1.2, \phantom{+}0.2)$ \\
6 & $(\phantom{+}2.2, -1.8)$ \\
7 & $(\phantom{+}3.0, \phantom{+}1.4)$ \\
8 & $(\phantom{+}0.9, \phantom{+}3.2)$ \\
\bottomrule
\end{tabular}
\caption{Generator means. The geometry is shown in Figure~\ref{fig:geometry}.}
\label{tab:means}
\end{table}

\subsection{Objectives}

\textbf{RKE (target-free diversity).} The loss is the pure quadratic form
\[
L_{\mathrm{RKE}}(\alpha) = \alpha^\top K \alpha,
\]
where $K_{ij} = \E_{X\sim P_i,\,X'\sim P_j}[\kappa(X,X')]$ with Gaussian RBF kernel $\kappa(x,x') = \exp\bigl(-\frac{\|x-x'\|^2}{2h^2}\bigr)$, $h=1.25$. The true kernel matrix is computed in closed form via
\[
K_{ij} = \Bigl(\frac{h^2}{h^2+2\sigma_g^2}\Bigr)^{d/2} \exp\!\Bigl(-\frac{\|\mu_i-\mu_j\|^2}{2(h^2+2\sigma_g^2)}\Bigr).
\]
The linear term is $f \equiv 0$. The range parameters are $\Delta_\kappa=1$, $\Delta_f=0$, $\Delta_L=2$.

\textbf{MMD to a target mixture.} The loss is the squared maximum mean discrepancy to a fixed target distribution $Q$, which is a 3-component Gaussian mixture:
\[
Q = 0.45\,P_3 + 0.35\,P_5 + 0.20\,P_7,
\qquad (\sigma_q=0.40).
\]
The quadratic form is
\[
L_{\mathrm{MMD}}(\alpha) = \alpha^\top K\alpha - 2b^\top\alpha + c,
\]
where $b_i = \E_{X\sim P_i,\,Y\sim Q}[\kappa(X,Y)]$ and $c = \E_{Y,Y'\sim Q}[\kappa(Y,Y')]$. The range parameters are $\Delta_\kappa=1$, $\Delta_f=2$, $\Delta_L=4$.

Crucially, under the factorization $K = M^\top M$ and the representation $y = M\alpha_s^\star$, this objective takes the least-squares form
\[
L_{\mathrm{MMD}}(\alpha) = \|M\alpha - y\|_2^2 - \|y\|_2^2,
\]
with $\alpha_s^\star = (0,0,0.45,0,0.35,0,0.20,0)$ achieving zero loss ($L(\alpha_s^\star) \approx 0$). The oracle is exactly $s=3$-sparse and the sparse and full optima coincide.

\subsection{Oracles}

Table~\ref{tab:oracles} summarizes the true optimum information. The sparse oracle is obtained by solving $\min_{\alpha\in\Delta_{m,s}} L(\alpha)$ exactly via support enumeration (92 supports checked for $m=8$, $s=3$). The full oracle solves $\min_{\alpha\in\Delta_m} L(\alpha)$.

\begin{table}[H]
\centering
\begin{tabular}{l|c|c}
\toprule
 & RKE & MMD \\
\midrule
$L(\alpha_s^\star)$ & $0.2770$ & $\approx 0$ \\
$\alpha_s^\star$ (support) & $\{1,6,8\}$ & $\{3,5,7\}$ \\
$\alpha_s^\star$ (weights) & $(0.33,0.33,0.33)$ & $(0.45,0.35,0.20)$ \\
$L(\alpha^\star)$ & $0.1655$ & $\approx 0$ \\
Best standalone arm & arm 1 & arm 3 \\
Standalone regret & $0.5530$ & $0.3532$ \\
\midrule
$\Delta_\kappa$ & $1.0$ & $1.0$ \\
$\Delta_f$ & $0.0$ & $2.0$ \\
$\Delta_L$ & $2.0$ & $4.0$ \\
\bottomrule
\end{tabular}
\caption{True oracle information for both objectives.}
\label{tab:oracles}
\end{table}

\subsection{Algorithms Compared}

Four methods are evaluated:

\begin{enumerate}[label=\textbf{A\arabic*}, leftmargin=*]
    \item \textbf{Fully-Corrective Sparse Mixture-UCB} (FC-Sparse, Algorithm~2). Greedily builds an active set of size $\le s$ using optimistic reduced-gradient selection and fully-corrective simplex QP corrections. Sampling distribution $q_{t+1} = \alpha_t$.
    \item \textbf{Coherence-Structured FC with Forced Exploration} (Coherence, Algorithm~3). Uses the \emph{unoptimistic} objective $\widehat L_t$ for NNOMP-style support discovery, then applies a final optimistic correction over the discovered face. Samples from $q_{t+1} = (1-\eta_t)\alpha_t + \eta_t \cdot \1/m$ with $\eta_t=0.05$ (constant forced exploration). Warm start: $n_0=2$ initial pulls per arm.
    \item \textbf{Non-sparse Mixture-UCB} (Full). Solves the optimistic QP over the full simplex $\Delta_m$. Serves as an upper baseline for what is achievable without sparsity constraints.
    \item \textbf{Best Standalone Generator} (Standalone). Deploys $e_{i^\star}$ where $i^\star = \argmin_i L(e_i)$ at every round.
\end{enumerate}

\subsection{Common Parameters}
\begin{itemize}[leftmargin=*]
    \item Horizon $T = 2000$, Monte Carlo repetitions $\text{reps} = 10$.
    \item Confidence parameter $\beta = 4.0$.
    \item Sparsity level $s = 3$.
    \item All methods share the same sample bank within each repetition for fair comparison.
\end{itemize}

\subsection{Regret Definition}
All methods are evaluated against the sparse benchmark:
\[
R_T(s) = \frac{1}{T}\sum_{r=1}^T \E[L(q_r)] - L(\alpha_s^\star),
\]
where $q_r$ is the deployed distribution at round $r$. The non-sparse baseline is measured against the \emph{same} sparse benchmark, so its regret can go below zero when it beats the sparse optimum.

\pagebreak

\section{Results}

\begin{figure}[H]
    \centering
    \begin{subfigure}[b]{0.48\textwidth}
        \centering
        \includegraphics[width=\textwidth]{results_full/rke_geometry.png}
        \caption{RKE}
        \label{fig:rke-geom}
    \end{subfigure}
    \hfill
    \begin{subfigure}[b]{0.48\textwidth}
        \centering
        \includegraphics[width=\textwidth]{results_full/mmd_geometry.png}
        \caption{MMD}
        \label{fig:mmd-geom}
    \end{subfigure}
    \caption{2D Gaussian generator geometry. Circles show one standard deviation. For MMD, starred points denote the three target components with radii proportional to their mixture weights $(0.45,0.35,0.20)$.}
    \label{fig:geometry}
\end{figure}

\subsection{Regret Curves}

\begin{figure}[H]
    \centering
    \begin{subfigure}[b]{0.48\textwidth}
        \centering
        \includegraphics[width=\textwidth]{results_full/rke_regret.png}
        \caption{RKE}
        \label{fig:rke-regret}
    \end{subfigure}
    \hfill
    \begin{subfigure}[b]{0.48\textwidth}
        \centering
        \includegraphics[width=\textwidth]{results_full/mmd_regret.png}
        \caption{MMD}
        \label{fig:mmd-regret}
    \end{subfigure}
    \caption{Regret curves $R_t = \frac{1}{t}\sum_{r=1}^t L(q_r) - L(\alpha_s^\star)$. Shaded bands show $\pm 1$ standard error over 10 repetitions. The dashed horizontal line is the sparse oracle benchmark.}
    \label{fig:regret}
\end{figure}

Table~\ref{tab:final-regret} reports the final mean regret at $T=2000$ with standard errors.

\begin{table}[H]
\centering
\begin{tabular}{l|c|c}
\toprule
\textbf{Method} & \textbf{RKE} & \textbf{MMD} \\
\midrule
FC-Sparse & $0.0176 \pm 0.0002$ & $0.2004 \pm 0.0017$ \\
Non-sparse (Full) & $-0.0994 \pm 0.0004$ & $0.1094 \pm 0.0013$ \\
\textbf{Coherence (Alg.~3)} & $\mathbf{0.0217 \pm 0.0015}$ & $\mathbf{0.0142 \pm 0.0005}$ \\
Standalone & $0.5530 \pm 0.0000$ & $0.3532 \pm 0.0000$ \\
\bottomrule
\end{tabular}
\caption{Final mean regret $\pm$ standard error at $T=2000$, 10 repetitions.}
\label{tab:final-regret}
\end{table}

\subsection{MMD: Why Coherence Excels}

On MMD, the coherence-structured algorithm achieves a dramatic 7--14$\times$ improvement over the next-best sparse method. Figure~\ref{fig:mmd-mixtures} reveals the mechanism.

\begin{figure}[H]
    \centering
    \begin{subfigure}[b]{0.30\textwidth}
        \centering
        \includegraphics[width=\textwidth]{results_full/mmd_oracle_sparse.png}
        \caption{True sparse oracle}
    \end{subfigure}
    \hfill
    \begin{subfigure}[b]{0.30\textwidth}
        \centering
        \includegraphics[width=\textwidth]{results_full/mmd_final_mixtures.png}
        \caption{Coherence final mixture}
    \end{subfigure}
    \caption{MMD: (a) The true sparse oracle supported on generators $\{3,5,7\}$ with weights $(0.45,0.35,0.20)$. (b) The coherence-structured algorithm's average final mixture, which recovers exactly this support.}
    \label{fig:mmd-mixtures}
\end{figure}

\begin{table}[H]
\centering
\begin{tabular}{l|c|c}
\toprule
\textbf{Method} & \textbf{Final weights (arms 1--8)} & \textbf{Support arms \% of pulls} \\
\midrule
FC-Sparse & $(.08,.14,.27,.08,.32,0,.06,.05)$ & 58\% \\
Full & $(.09,.08,.25,.06,.21,.09,.14,.07)$ & 55\% \\
\textbf{Coherence} & $\mathbf{(0,0,.42,0,.34,0,.25,0)}$ & \textbf{96\%} \\
\bottomrule
\end{tabular}
\caption{Average final deployed mixtures and pull allocation on MMD. Coherence discovers the true support $\{3,5,7\}$ and concentrates 96\% of its 2000 pulls there, giving much tighter empirical estimates of $K$ and $f$ on the relevant subspace.}
\label{tab:mmd-weights}
\end{table}

The coherence algorithm succeeds because:
\begin{enumerate}[leftmargin=*]
    \item The MMD objective factorizes as $\|M\alpha - y\|_2^2$ with $y = M\alpha_s^\star$, satisfying the exact nonnegative sparse representation condition.
    \item The NNOMP-style unoptimistic gradient rule $\argmin_j [\nabla \widehat L_t(\widehat b_U)]_j$ correctly identifies the three target generators $\{3,5,7\}$.
    \item Once the support is recovered, the final optimistic correction over $\Delta(\{3,5,7\})$ is exact---the regret proof reduces to that of Exact Sparse Mixture-UCB, with no $O(\Delta_\kappa/s)$ structural term.
    \item Forced exploration ($\eta = 0.05$) ensures all arms are sufficiently observed for score concentration, while 95\% of pulls go to the three support arms, yielding far tighter concentration than the uniform exploration of standard Mixture-UCB.
\end{enumerate}

\subsection{RKE: When Coherence Offers No Advantage}

On RKE, coherence is comparable to FC-Sparse ($0.022$ vs.\ $0.018$). The RKE objective lacks the least-squares structure:
\begin{itemize}[leftmargin=*]
    \item $f \equiv 0$, so there is no ``target $y$'' for the NNOMP residual-correlation rule to track.
    \item Gradient components $\nabla L(e_i) = 2K_{ii}$ differ only through the diagonal of $K$, which is nearly constant ($K_{ii} \approx 0.83$) because the RBF kernel is translation-invariant.
    \item There is no coherence-based support-recovery guarantee---the NNOMP path is essentially unguided.
    \item The fully-corrective mechanism still yields a valid $s$-sparse solution, but there is no structural advantage over the standard optimistic FC.
\end{itemize}

This is consistent with the theory: the coherence-structured theorem requires $K=M^\top M$, a least-squares-compatible linear term $f = -2M^\top y$, and an exact nonnegative sparse representation $y = M\alpha_s^\star$. RKE satisfies the factorization but not the linear-term compatibility.

\section{Computational Cost}

Table~\ref{tab:compute} reports wall-clock times per repetition. The full-simplex method uses active-set enumeration (255 subsets for $m=8$) at every round; FC-Sparse and Coherence solve small simplex QPs on active sets of size $\le 3$ (at most 7 subsets per correction). All experiments used a single CPU core (Apple M-series).

\begin{table}[H]
\centering
\begin{tabular}{l|c}
\toprule
\textbf{Method} & \textbf{Seconds per rep} ($T=2000$) \\
\midrule
FC-Sparse & $\sim 20$ \\
Full & $\sim 20$ \\
Coherence & $\sim 20$ \\
\bottomrule
\end{tabular}
\caption{Wall-clock time per repetition (10 reps averaged). All methods have comparable cost since the bottleneck is the active-set simplex QP solver used for the full-simplex method ($m=8$, 255 subsets per round), not the support discovery overhead.}
\label{tab:compute}
\end{table}

\section{Summary}

The coherence-structured forced-exploration algorithm (Algorithm~3) fulfills its theoretical promise on the MMD objective: by leveraging the least-squares factorization and exact sparse representation, it recovers the oracle support and concentrates samples there, eliminating the $O(\Delta_\kappa/s)$ Frank-Wolfe structural term. The resulting regret ($0.0142$) is $14\times$ lower than standard FC-Sparse ($0.2004$) and even $7.7\times$ lower than the non-sparse baseline ($0.1094$).

On the RKE objective, which lacks the required structural assumptions, coherence performs comparably to FC-Sparse, confirming that the method degrades gracefully when the coherence conditions are not met.

A natural next experiment is to sweep multiple sparsity levels $s \in \{1,\dots,5\}$ and candidate-set sizes $m \in \{8,12,16\}$ to characterize how the coherence advantage scales with problem dimension.

\end{document}
